\section*{Chapitre \ref{ch:s2:discretionary-macro} --- Macro discrétionnaire}

\begin{solution}{exo:s2:discretionary-macro:1}
$99\sqrt{0.5}/\sqrt{2\pi} = 27.9$ points de base.
\end{solution}

\begin{solution}{exo:s2:discretionary-macro:2}
La position avec le stop à 25 points de base est deux fois plus grande que celle avec le stop à 50 points de base~: chacune perd une unité lorsqu'elle est stoppée.
\end{solution}

\begin{solution}{exo:s2:discretionary-macro:3}
Le spread coûte sa pleine largeur, soit 50 points de base~: $50/18.2 - 1 = 1.75$ unités de budget, nettes de la prime.
\end{solution}

\begin{solution}{exo:s2:discretionary-macro:4}
La probabilité qu'un chemin touche pour la première fois un niveau décroît rapidement avec la distance de ce niveau, mesurée en écart type du chemin~: à six mois, l'écart type
du taux est d'environ 64 points de base, de sorte qu'un stop à 10 points de base est touché sur la plupart des chemins et un stop à 100 points de base presque jamais.
\end{solution}

\begin{solution}{exo:s2:discretionary-macro:5}
Il coûte moins cher (17.2 contre 27.9 points de base), de sorte que le budget en achète davantage et qu'une forte baisse rapporte davantage par unité~; mais le taux doit baisser
de plus de 25 points de base pour qu'il rapporte quoi que ce soit, de sorte qu'il expire sans valeur plus souvent.
\end{solution}

\begin{solution}{exo:s2:discretionary-macro:6}
Sur six mois, le bruit propre du taux (environ 64 points de base d'écart type) est plus grand que la baisse de 50 points de base attendue~; sur environ 21~\%
des chemins, le bruit l'emporte donc, même lorsque la dérive est exactement celle prévue par la vue.
\end{solution}

\begin{solution}{exo:s2:discretionary-macro:7}
Sans dispersion, le taux finit plus bas sur 78.6~\% des chemins au lieu de 73.4~\%, mais la plupart des gains par unité de budget baissent~: les contrats à terme à stop serré passent
de 1.68 à 1.45, l'option à la monnaie de 1.26 à 1.07, l'option hors de la monnaie de 1.69 à 1.32, les contrats à terme à stop de 50 points de base de 0.98 à 0.91~; seul
le spread plafonné monte, de 0.71 à 0.79. La dispersion de la dérive crée les grands mouvements que paient le plus les expressions convexes~; la précision aide
les expressions qui abandonnent les grands mouvements.
\end{solution}

\begin{solution}{exo:s2:discretionary-macro:8}
Avoir raison sur la direction n'est pas la même chose qu'avoir raison sur l'ampleur, le moment et le chemin~; et une année de résultats est trop peu pour distinguer
le talent du bruit. Le journal des transactions doit montrer si c'est l'expression, la taille ou le stop qui a transformé une vue juste en perte.
\end{solution}

\begin{solution}{pb:s2:discretionary-macro:1}
\begin{enumerate}
\item L'instrument, le strike, l'échéance et la taille qui transforment une vue en position.
\item Contrats à terme~: la taille~; options~: la taille avec une perte plafonnée~; stops~: le chemin~; spreads~: un mouvement jusqu'à un point.
\item Une baisse de 50 points de base en six mois, une dispersion de la dérive de 50, une volatilité de 90 par an.
\item Sur 73.4~\% des chemins.
\item La prime que l'on est prêt à dépenser pour les options.
\item Les options par leur prime~; les contrats à terme stoppés par leur stop~; les contrats à terme non stoppés par une perte de deux écarts types.
\item 27.9, 17.2 et 18.2 points de base.
\item Chaque définition du risque dimensionne les positions différemment.
\item Une règle de réduction d'une position, fixée à l'avance et respectée.
\item 55.6~\%, 26.5~\%, 5.6~\%, 0.7~\% et 0.1~\% pour 10, 25, 50, 75 et 100 points de base.
\item Un cadre pour la valeur des règles de stop-loss~; certaines augmentent le rendement espéré tout en réduisant la volatilité aux fréquences plus longues.
\item Un stop plus serré signifie une position plus grande pour le même budget.
\item \textbf{Résultat nommé.} Par unité de budget de risque, les contrats à terme rapportent 0.39, les contrats à terme avec un stop de 25 ou 50 points de base 1.68 et 0.98, les options 1.26 (à la
monnaie) et 1.69 (hors de la monnaie) et le spread 0.71~; un stop de 25 points de base coupe court 26.5~\% des vues justes, un stop de 50 points de base 5.6~\%.
\item Taille~: contrats à terme ou options~; chemin~: options plutôt que stops serrés.
\item La vue a raison plus souvent~; les expressions convexes rapportent moins.
\item La vue, son horizon et sa confiance, l'expression et sa raison, et le résultat de chacune.
\item Séparer le taux de réussite des vues du rendement des expressions, sur de nombreuses transactions.
\item Il s'agit d'exprimer n'importe quelle vue, et non une stratégie dotée d'une source de rendement.
\item Le chapitre~16 remplace les vues du gérant par des règles.
\item L'expression, la taille et le stop.
\end{enumerate}
\end{solution}

\begin{solution}{iq:s2:discretionary-macro:1}
Recevoir le fixe dans les swaps à court terme ou acheter les contrats à terme pertinents (par exemple, les contrats à terme SOFR pour les dates de réunion), ou acheter des options de réception ou des calls sur
les contrats à terme~; on choisit selon le degré de certitude sur le moment et selon la perte que l'on peut supporter.
\end{solution}

\begin{solution}{iq:s2:discretionary-macro:2}
Lorsque le chemin est incertain (le marché peut d'abord aller contre la vue), lorsque la perte doit être plafonnée à l'avance, ou lorsque la vue porte sur un grand
mouvement dont l'ampleur est incertaine.
\end{solution}

\begin{solution}{iq:s2:discretionary-macro:3}
À partir du bruit du chemin attendu sur l'horizon~: assez large pour que le bruit les déclenche rarement, dimensionné de sorte qu'un stop déclenché coûte une part fixée de la
limite de perte du livre~; on les revoit à la lumière de l'historique des vues stoppées.
\end{solution}

\begin{solution}{iq:s2:discretionary-macro:4}
Consigner les vues au moment où elles sont formées et les évaluer par rapport aux résultats sur de nombreuses transactions, relativement à un taux de base et à une expression aléatoire~; séparer le taux de réussite du
gain par transaction.
\end{solution}

\begin{solution}{iq:s2:discretionary-macro:5}
Chaque entrée~: l'heure, la vue (direction, taille, horizon, confiance), l'expression et sa taille, le stop et l'objectif, puis les sorties avec leur raison et le P\&L~; étiquetées pour que les vues
et les expressions puissent être évaluées séparément.
\end{solution}

\begin{solution}{iq:s2:discretionary-macro:6}
Par le principe de réflexion, $P(\max_{t \le T} W_t \ge s) = 2\,(1 - N(s/(\sigma\sqrt T)))$. Elle décroît vite dès que $s$ dépasse un écart type
$\sigma\sqrt T$ (environ 64 points de base ici), ce qui explique pourquoi la courbe de sortie par stop chute fortement entre 10 et 50 points de base~; la dérive de la vue la fait
décroître plus vite encore.
\end{solution}
